\documentclass[10pt]{article}
\usepackage[margin=1in]{geometry}
\usepackage{booktabs}
\usepackage{graphicx}
\usepackage{amsmath}
\begin{document}

\begin{center}
{\large\bf Supplementary Material: GA-DCHG Implementation Details and Additional Results}
\end{center}

\vspace{1em}

\noindent\textbf{1 Implementation details}

All models are trained with the Adam optimizer at learning rate $2\times10^{-3}$ for 60 epochs with batch size 16; the 1-Wasserstein alignment term is active from the first epoch. Every protocol runs over three seeds $\{0,1,2\}$, and the checkpoint is selected by source validation accuracy and evaluated on the target test set. Hypergraph caches are precomputed once per protocol; all reported numbers are produced by the released pipeline from the raw CWRU recordings.

\vspace{0.5em}

\noindent\textbf{1.1 Multi-source CWRU screening (cwru\_0123)}

\begin{table}[h]
\centering
\begin{tabular}{lcc}
\toprule
Descriptor & W1 distance & Anchor \\
\midrule
$\kappa$ & 1.230 & \\
$\chi$   & 1.097 & \\
$\rho$   & 0.848 & anchor \\
\bottomrule
\end{tabular}
\end{table}

\vspace{0.5em}

\noindent\textbf{1.2 Multi-source CWRU cross-domain evaluation (cwru\_0123)}

\begin{table}[h]
\centering
\begin{tabular}{lcc}
\toprule
Method & Acc. & Gap (vs 1D-CNN) \\
\midrule
1D-CNN            & 0.788 & $+0.000$ \\
DynHGNN+          & 0.862 & $+0.074$ \\
GA-DCHG           & \textbf{0.902} & $\mathbf{+0.114}$ \\
\hspace{1em}$-$align      & 0.892 & $+0.104$ \\
\hspace{1em}$-$temporal   & 0.872 & $+0.084$ \\
\hspace{1em}$-$anchor     & 0.848 & $+0.060$ \\
\bottomrule
\end{tabular}
\end{table}

\vspace{0.5em}

\noindent\textbf{2 Notes}

The multi-source CWRU protocol (loads 0,1,2 $\to$ 3) confirms the main findings on a second CWRU shift: the structural baseline and GA-DCHG again exceed the 1D-CNN by a clear margin, and GA-DCHG is the best model. Removing the data-selected anchor ($-$anchor, 0.848) or temporal propagation ($-$temporal, 0.872) produces the largest drops, consistent with the main text's claim that the anchor screening and temporal propagation carry part of the cross-domain benefit.

\end{document}
